\section{Правила оформления}

\label{NB-rules}

{ \color{red}

Напоминаю короткие правила.

1) Местоимение "Я" запрещено - правильно "МЫ":

\textbf{мы показали, что... }

2) Категорически запрещено использовать сленг и сокращения кроме общепринятых. 

3) Грамматику и пунктуацию проверяем сразу а не "потом"



Графики, диаграммы, блок - схемы загружаются в векторном формате EPS или PDF. Сами файлы должны лежать в папке eps  или pdf.

Растровые рисунки в форматах png и jpg должны лежать в каталоге figures.

На рисунках должны отсутствовать подпись рисунка и подписи осей на графиках (как это модно в NumPy). 

В NumPy графиках оси черные 2pt - сетка серая 1pt.

199\,129 вот так отделяются группы по три цифры в длинных числах

Формулы оформляем окружением equation.

Формулы в составе строки включаем, как \$formula\$.

Ссылки на литературу вставляем строго через bib файл.

Используемые кавычки Только 
``ТАКИЕ''



}